%! Convolutional Neural Nets — Unit 8, Lesson 8.2 — Homework (Answer Key)
\documentclass[10pt]{article}
\usepackage{linalg-article}
\usepackage{linalg-key}   % pulls in -boxes; do NOT also load -boxes
\newcommand{\TallMath}[1]{$\displaystyle #1\rule[-1.4em]{0pt}{3.2em}$}
\newcommand{\vv}{\mathbf{v}}
\newcommand{\bb}{\mathbf{b}}
\newcommand{\ww}{\mathbf{w}}
\newcommand{\zero}{\mathbf{0}}
\newcommand{\relu}{\mathrm{ReLU}}

\begin{document}

\pageheader{Unit 8, Lesson 8.2}{Homework --- Answer Key}
\namedateperiod

\vspace{0.12in}

\begin{notesbox}{Practice}
A \textbf{convolution} slides a small \textbf{filter} across a signal, taking a dot product at each spot. Written
as a matrix, the filter's weights \textbf{repeat down the diagonals} --- the same few weights reused everywhere
(\textbf{weight sharing}). A convolutional layer is $\relu(A\vv+\bb)$ with $A$ a convolution matrix.

\begin{enumerate}[leftmargin=1.9em, itemsep=11pt]
  \item \textbf{Slide the edge filter.} Use $\ww=[-1,1]$ (output $=-v_i+v_{i+1}$). Give the output and say
        where it fires (if anywhere).
    \begin{enumerate}[label=(\alph*), leftmargin=1.8em, itemsep=4pt]
      \item $\vv=[4,4,4,4]$ \hfill \ans{output $=[0,0,0]$; never fires (the signal is flat)}
      \item $\vv=[1,1,6,6]$ \hfill \ans{output $=[0,5,0]$; fires at spot $2$ (the edge $1\to6$)}
    \end{enumerate}
  \item \textbf{Smoothing filter.} Slide the averaging filter $\ww=\tfrac12[1,1]$
        (output $=\tfrac{v_i+v_{i+1}}{2}$) across $\vv=[2,6,6,2]$. Give the output. How is its effect
        different from the edge filter's? \ansline{Output $[\,4,\,6,\,4\,]$. The averaging filter \emph{smooths} (blurs) the signal --- it replaces sharp changes with gentler ones --- whereas the edge filter \emph{sharpens}, firing only at jumps.}
  \item \textbf{Write it as a matrix.} For the filter $\ww=[-1,1]$ and a length-$4$ signal, write the
        $3\times4$ convolution matrix $A$ (the two numbers slide over one step each row). Then compute
        $A\vv$ for $\vv=[1,1,6,6]$. What do you notice about the weights in $A$? \ansline{$A=\begin{bsmallmatrix} -1 & 1 & 0 & 0\\ 0 & -1 & 1 & 0\\ 0 & 0 & -1 & 1 \end{bsmallmatrix}$, and $A\vv=[0,5,0]$. The same two numbers $-1,1$ repeat down the diagonals (shifted one step per row) --- only $2$ distinct weights, shared across the whole matrix.}
  \item \textbf{Count the weights.} (a) For the $3\times4$ matrix in problem~3, how many separate weights
        would a \emph{full} layer have, and how many does the convolution use? (b) A small $28\times28$ image
        has $784$ pixels; a full layer to $784$ outputs uses $784^2=614{,}656$ weights. A $3\times3$ filter
        uses how many? \ansline{(a) Full: $3\times4=12$ weights; convolution: $2$. (b) A $3\times3$ filter uses $9$ weights --- versus $614{,}656$ for the full layer.}
  \item \textbf{Explain.} (a) Why does reusing one filter let a network detect the same pattern \emph{wherever}
        it appears in the image? (b) In the layer $\relu(A\vv+\bb)$, how does a negative bias act like a
        \emph{threshold} on the detector? \ansline{(a) The identical filter weights are applied at every position, so a pattern produces the same response no matter where it sits --- one detector covers all locations (translation invariance). (b) The bias subtracts a constant before the ReLU, so only responses \emph{above} that amount stay positive and survive; weaker ones become negative and are zeroed --- the bias sets how strong the pattern must be to ``count.''}
\end{enumerate}
\end{notesbox}

\begin{extensionbox}
\textbf{One filter finds one pattern.}
\begin{enumerate}[label=(\alph*), leftmargin=1.8em, itemsep=5pt]
  \item A $5\times5$ filter has how many weights? Does that number grow when the image gets bigger? Why does
        that matter for scanning a huge photo?
  \item Apply the full detector $\relu(A\vv+\bb)$ with the edge filter $\ww=[-1,1]$ and bias $\bb=-3$ to
        $\vv=[2,2,2,7,7]$. (First convolve, then add the bias, then ReLU.) Which spot survives, and what does
        the surviving value tell you?
\end{enumerate}
\ansline{(a) $5\times5=25$ weights. It does \emph{not} grow with image size --- the same $25$ weights are reused everywhere --- so one small filter can scan an image of any size cheaply.}
\ansline{(b) Convolve: $[2,2,2,7,7]\to[0,0,5,0]$. Add bias $-3$: $[-3,-3,2,-3]$. ReLU: $[0,0,2,0]$. Only the third spot survives; the value $2$ says the edge there was strong enough to clear the threshold ($5>3$), so the detector reports the pattern at that location.}
\end{extensionbox}

\begin{spiralbox}
\textbf{Coming next --- Lesson 8.3: Minimizing Loss by Gradient Descent.} We have built the pieces (layers,
folds, filters); next we \emph{choose} the weights --- rolling ``downhill'' on the loss to make the network
fit the data.
\end{spiralbox}

\begin{teachernote}
Item~1 drills the slide-and-fire mechanic (flat $\to$ all zeros; a jump $\to$ a spike at the edge). Item~2
contrasts the smoothing filter (blurs) with the edge filter (sharpens) --- same slide, opposite jobs. Item~3
is the linear-algebra core: the convolution \emph{is} a matrix, but a constant-diagonal one with only $2$
shared weights. Item~4 is the parameter-count payoff ($12$ vs.\ $2$; $614{,}656$ vs.\ $9$). Item~5 collects
the two big ideas: weight sharing $=$ detect-anywhere, and the bias $=$ a detection threshold under the ReLU.
The extension reinforces that filter size is fixed regardless of image size, and walks the full
$\relu(A\vv+\bb)$ pipeline once more. Most common slips: sign/order in $-v_i+v_{i+1}$ (``next minus
current''); writing a full matrix instead of the repeated diagonal; applying the bias \emph{after} the ReLU
instead of before.
\end{teachernote}

\end{document}
